% --- 补遗：走廊细节、层约束、天线与冗余 via ---

\subsection{定义与使用 Routing Corridor（详述）}
\subsubsection*{Defining and Using Routing Corridors}

Routing corridor 将特定网络的全局布线限制在一组相连矩形所定义的区域内，并可对每个矩形指定最小/最大布线层。主要用于在信号布线前布关键网。若 routing guide 与 corridor 冲突，corridor 优先。

\begin{lstlisting}
icc2_shell> create_routing_corridor -name corridor_a \
   -boundary { {10 10} {20 35} } \
   -min_layer_name M2 -max_layer_name M4
icc2_shell> create_routing_corridor_shape -routing_corridor corridor_a \
   -boundary { {20 25} {40 35} } \
   -min_layer_name M2 -max_layer_name M4
icc2_shell> add_to_routing_corridor corridor_a [get_nets {n1 n2}]
icc2_shell> check_routing_corridors RC_0
\end{lstlisting}

也可按路径创建：\opt{-path} + \opt{-width}；端帽样式 \opt{-start\_endcap}/\opt{-end\_endcap}（\cmd{flush}/\cmd{full\_width}/\cmd{half\_width}）。一个 net/supernet 只能属于一个 corridor，且 corridor 须覆盖其全部 pin。

修改：\cmd{create\_routing\_corridor\_shape}、\cmd{remove\_routing\_corridor\_shapes}、\cmd{add\_to\_routing\_corridor}、\cmd{remove\_from\_routing\_corridor}。报告：\cmd{report\_routing\_corridors}（\opt{-output} 可写出 Tcl）、\cmd{get\_routing\_corridors -of\_objects}。移除：\cmd{remove\_routing\_corridors -all}。

\subsection{布线层资源与冻结层（章内相关）}
\subsubsection*{Routing Layer Resources and Freeze Layers}

全局/网络/时钟的最小最大布线层分别由 \cmd{set\_ignored\_layers}、\cmd{set\_routing\_rule} 与 \cmd{set\_clock\_routing\_rules} 的 \opt{-min\_routing\_layer}/\opt{-max\_routing\_layer} 设置；软/硬模式由 \appopt{route.common.global\_min\_layer\_mode}、\appopt{route.common.net\_min\_layer\_mode} 等控制。冻结层由 \appopt{route.common.freeze\_layer\_by\_layer\_name} 与 \appopt{route.common.freeze\_via\_to\_frozen\_layer\_by\_layer\_name} 指定。电压域穿越与层优选方向的详细设置见第~2~章相关主题。

\subsection{天线规则与冗余 Via（章内要点）}
\subsubsection*{Antenna Rules and Redundant Vias}

详细布线阶段并发处理天线规则，并在报告中给出天线违例与冗余 via 转换率。用 \cmd{check\_routes} 的 \opt{-antenna} 可启用/关闭天线检查。冗余 via 可在信号布线期间并发插入；详细布线结束报告含 ``Redundant via conversion report''（按层与 weight 的优化转换率）。更完整的天线修复与冗余 via 专项流程见原书交叉引用主题及第~6~章芯片收尾相关内容。
